% !TEX root = ../main.tex
%\vspace{-2mm}
\section{Introduction}
\vspace{-3mm}
  \begin{figure}[ht!]
    \centering
    \includegraphics[width=.7\linewidth, trim = 0.9cm 0.5cm 4.4cm 0.4cm, clip]{images/teaser.pdf}
    \caption{ DLMs require many reverse-diffusion steps at inference. \emph{IDLM} trains a few-step generator via inverse distillation to match the teacher’s generation quality with substantially fewer steps.}
    \vspace{-4.5mm}
  \label{fig:teaser}
  \end{figure}


Generative modeling for discrete data, such as natural language, has seen widespread adoption, driven by the rapid advancement of large language models~\citep{touvron2023llama, groeneveld2024olmo, lozhkov2024starcoder}. 

Recently, Diffusion Language Models 
(DLMs; \citealp{sohl2015deep, austin2021structured, campbell2022continuous, lou2024discrete}) have emerged as a promising framework for modeling the distribution of the text data. A general formulation for a broad class of DLMs was introduced by~\citep{lou2024discrete}. However, the general formulation is often intractable in practice and introduces significant theoretical complexity. To address this, subsequent work~\citep{sahoo2024simple, sahoo2025the, schiff2024discreteguidance, shi2024simplified} has focused on specific types of diffusion processes, enabling simplified theoretical analyses and more tractable implementations. As a result, DLMs have achieved competitive performance in text generation, comparable to that of autoregressive language models~\citep{nie2025large}. However, the inherently iterative nature of reverse diffusion requires hundreds or even thousands of sampling steps at inference time, resulting in considerable latency that limits their practicality.

On the other hand, in the continuous domain, the issue of slow inference in diffusion models~\citep{ho2020denoising, song2020score, karras2022elucidating} has been the focus of extensive research~\citep{song2020denoising, luhman2021knowledge, zhang2022fast, lu2022dpm}. Among the various approaches, \emph{distillation} has emerged as one of the most effective solutions~\citep{song2023consistency, salimans2022progressive, luhman2021knowledge, berthelot2023tract, xie2024distillation}. However, directly applying such techniques in the discrete domain is challenging due to the non-differentiable nature of discrete data. Recent efforts have begun to address this limitation for DLMs. Notably, SDTT~\citep{deschenaux2024beyond} and Duo-DCD~\citep{sahoo2025the} propose adaptations of consistency-based distillation methods~\citep{song2023consistency, song2023improved, kim2023consistency} to the discrete setting, demonstrating initial progress toward enabling fast inference in DLMs.

  In this work, we investigate an alternative distillation paradigm by extending the theoretical framework of \emph{Inverse Distillation}~\citep{gushchin2025inverse, kornilov2025universal}, originally developed for continuous diffusion models, to the discrete domain. However, this direct extension introduces several nontrivial challenges arising from the intrinsic properties of discrete data. First, \underline{theoretical concerns} emerge regarding the validity of the training objective, particularly due to the potential non-uniqueness of its minimizers. Second, \underline{practical challenges} arise from the non-smoothness of the discrete space, which complicates the optimization. 
  
  \textbf{Contributions.} To address the above-mentioned issues, we present \emph{Inverse-distilled Diffusion Language Models (IDLM)}, a principled framework for accelerating a broad class of discrete diffusion models. In summary, our core contributions are as follows:

  \vspace{-0.8em}

  \begin{enumerate}
    \setlength{\itemsep}{1pt}
    \setlength{\parskip}{1pt}
    \item We theoretically establish the uniqueness of the global optimum for the IDLM objective, validating the correctness of the proposed optimization procedure (\wasyparagraph\ref{sec:theoretical extension idlm}).

    \item We propose gradient-stable relaxations that enable effective training through (i) a \emph{simplex-parameterized} generator and (ii) a principled mechanism for \emph{backpropagating gradients through the discrete forward process} (\wasyparagraph\ref{sec:simplex relaxation}).
    \item We empirically demonstrate that IDLM achieves competitive \emph{generation quality}, matching the performance of 1024-step teacher models while requiring up to $64\times$ fewer sampling steps (\wasyparagraph\ref{sec:exps}).
  \end{enumerate}
